\documentclass[12pt]{article}
\usepackage{amsmath,amssymb,amsfonts}
\usepackage[margin=1in]{geometry}

\begin{document}

\textbf{Chapter 6, Problem 6.} Show that if we write $c = re^{i\theta}$, the Cauchy--Riemann equations become, in terms of $r$ and $\theta$,
\[
\frac{\partial u}{\partial r} = \frac{1}{r}\frac{\partial v}{\partial \theta}, \qquad \frac{\partial v}{\partial r} = -\frac{1}{r}\frac{\partial u}{\partial \theta}.
\]

\bigskip
\textbf{Solution.}

With $z = x + iy = re^{i\theta}$, we have $x = r\cos\theta$ and $y = r\sin\theta$.

The Cauchy--Riemann equations in Cartesian form are:
\[
\frac{\partial u}{\partial x} = \frac{\partial v}{\partial y}, \qquad \frac{\partial u}{\partial y} = -\frac{\partial v}{\partial x}.
\]

Using the chain rule:
\[
\frac{\partial u}{\partial r} = \frac{\partial u}{\partial x}\cos\theta + \frac{\partial u}{\partial y}\sin\theta
\]
\[
\frac{\partial u}{\partial \theta} = -\frac{\partial u}{\partial x}\,r\sin\theta + \frac{\partial u}{\partial y}\,r\cos\theta
\]

Similarly for $v$:
\[
\frac{\partial v}{\partial r} = \frac{\partial v}{\partial x}\cos\theta + \frac{\partial v}{\partial y}\sin\theta
\]
\[
\frac{\partial v}{\partial \theta} = -\frac{\partial v}{\partial x}\,r\sin\theta + \frac{\partial v}{\partial y}\,r\cos\theta
\]

Now apply the Cauchy--Riemann equations ($u_x = v_y$, $u_y = -v_x$):
\[
\frac{1}{r}\frac{\partial v}{\partial \theta} = \frac{1}{r}\left(-\frac{\partial v}{\partial x}\,r\sin\theta + \frac{\partial v}{\partial y}\,r\cos\theta\right) = -v_x\sin\theta + v_y\cos\theta
\]
\[
= u_y\sin\theta + u_x\cos\theta = \frac{\partial u}{\partial r}. \quad \checkmark
\]

Similarly:
\[
-\frac{1}{r}\frac{\partial u}{\partial \theta} = -\frac{1}{r}\left(-u_x\,r\sin\theta + u_y\,r\cos\theta\right) = u_x\sin\theta - u_y\cos\theta
\]
\[
= v_y\sin\theta + v_x\cos\theta = \frac{\partial v}{\partial r}. \quad \checkmark
\]

This establishes the Cauchy--Riemann equations in polar form.

\end{document}
